\section{Growth with explicit arithmetic error bounds}
\label{sec:growth}

The aim of this section is stronger than a statement that the usual
asymptotic argument is effective.  We give inequalities with explicit
errors, and, for the principal asymptotic equivalences, an explicit
primitive recursive modulus.  The logarithms, exponentials, square roots,
integrals, and the function $W$ in the displayed formulas denote the
represented quantities described in the foundations section.  A
non-strict inequality between them means the corresponding family of
rational inequalities with arbitrarily small prescribed slack.  No
quantification over an unrestricted space of real functions is needed.

Throughout, for $n\geq3$, put
\[
 r=W(n),\qquad a=e^r=\frac nr,\qquad
 b=n(r+1),\qquad L=\log n.
\]
Thus $r>1$, $r+\log r=L$, and
\begin{equation}\label{eq:growth-r-bracket}
 L-\log L\leq r\leq L.
\end{equation}
The lower inequality follows from $\log r\leq\log L$.  The function
$W(n)$ can be computed without deciding equality between real numbers.
For example, invert $x\mapsto xe^x$ on the rational interval $[0,n]$.
If a current interval has width $w$ and midpoint $m$, approximate
$me^m-n$ to error at most $w/8$.  An approximation greater than $w/4$
or less than $-w/4$ certifies which half contains the root.  Otherwise
the root is within $3w/8$ of $m$, since the derivative of $xe^x$ is at
least one for $x\geq0$.  Intersect with that smaller interval.  Each
step decreases the width by a factor at most $3/4$, so a predetermined
primitive recursive number of steps suffices for any requested error.
Existence here is the compatibility of these explicit brackets; it is
not an application of an unrestricted intermediate value principle.

\subsection{A logarithmic envelope requiring no saddle-point argument}

We first record the elementary factorial bounds
\begin{equation}\label{eq:elementary-factorial-log}
 m\log m-m+1\leq\log(m!)
 \leq m\log m-m+1+\log m\qquad(m\geq1).
\end{equation}
They follow by bounding the increasing function $\log x$ on each of
the intervals $[j,j+1]$ and adding the resulting finite inequalities.
The integral $\int_1^m\log x\,dx=m\log m-m+1$ is obtained from its
explicit antiderivative.  Equivalently, the same proof uses rational
Riemann sums and their explicit error bounds.

\begin{theorem}[An explicit logarithmic envelope]\label{thm:log-envelope}
For every integer $n\geq3$, let
\[
 H(n)=n\left(r-1+\frac1r\right),\qquad r=W(n).
\]
Then
\begin{equation}\label{eq:log-envelope}
 H(n)-r-4\leq\log B_n\leq H(n)+\log n.
\end{equation}
In particular,
\begin{equation}\label{eq:log-envelope-symmetric}
 \left|\frac{\log B_n}{n}-
 \left(r-1+\frac1r\right)\right|
 \leq\frac{\log n+4}{n}.
\end{equation}
\end{theorem}

\begin{proof}
All coefficients of the exponential generating function are
nonnegative.  Consequently, for every positive $x$,
\[
 B_n\leq \frac{n!}{x^n}\exp(e^x-1).
\]
Set $x=r$ and use the upper bound in
\eqref{eq:elementary-factorial-log}.  Because
$\log(n/r)=r$ and $a=n/r$, this gives
\[
 \log B_n\leq n\log n-n+1+\log n+a-1-n\log r
 =H(n)+\log n.
\]

For the lower bound take one positive term of Dobi\'nski's formula.
There is a small constructive point in selecting that term: we do not
take the floor of an arbitrary represented real.  Compute a rational
lower approximation $u$ satisfying $a-1/4\leq u\leq a$, and take the
ordinary rational floor $k=\lfloor u\rfloor$.  Then
\[
 a-\frac54<k\leq a,\qquad k\geq\frac a2,\qquad k\geq1.
\]
Define
\[
 g(x)=n\log x-x\log x+x.
\]
We have $g'(a)=0$, $g(a)=H(n)$, and, for $k\leq x\leq a$,
\[
 |g''(x)|=\frac n{x^2}+\frac1x
 \leq\frac{6r}{a}.
\]
Taylor's formula with its integral remainder therefore gives
\[
 0\leq g(a)-g(k)
 \leq\frac{3r}{a}\left(\frac54\right)^2
 =\frac{75r}{16e^r}<2.
\]
Here $re^{-r}\leq e^{-1}$ for $r\geq1$, and $e>5/2$ suffices for
the last numerical bound.  Dobi\'nski and
\eqref{eq:elementary-factorial-log} now give
\[
 \log B_n\geq n\log k-\log(k!)-1
 \geq g(k)-\log k-2\geq H(n)-r-4.
\]
This proves both assertions.
\end{proof}

\subsection{The de Bruijn logarithmic expansion}

\begin{theorem}[The displayed de Bruijn expansion, with constants]
\label{thm:debruijn}
Let $L=\log n\geq16$ and $\ell=\log L$.  Then
\begin{align}
 \left|\frac{\log B_n}{n}
 -\left(L-\ell-1+\frac{\ell+1}{L}
             +\frac{\ell^2}{2L^2}\right)\right|
 &\leq\frac{6\ell^3}{L^3}+\frac{L+4}{n}
 \label{eq:debruijn-sharp-error}\\
 &\leq\frac{7\ell}{L^2}.
 \label{eq:debruijn-displayed-error}
\end{align}
Thus, with the usual notation interpreted by these bounds,
\[
 \frac{\log B_n}{n}
 =\log n-\log\log n-1+
 \frac{\log\log n+1}{\log n}
 +\frac12\left(\frac{\log\log n}{\log n}\right)^2
 +O\!\left(\frac{\log\log n}{(\log n)^2}\right).
\]
The integer threshold $n\geq2^{24}$ is sufficient.
\end{theorem}

\begin{proof}
We give the error calculation, so the order symbol conceals no
ineffective threshold.  Set
\[
 r_0=L-\ell+\frac\ell L+
             \frac{\ell^2/2-\ell}{L^2},\qquad
 u=\frac{r_0-L}{L}.
\]
For $L\geq16$ one has $\ell\geq2$ and $0\leq-u\leq\ell/L\leq1/4$.
The convergent geometric-series integral for $\log(1+u)$ gives
\[
 \left|\log(1+u)-u+\frac{u^2}{2}\right|
 \leq\frac{|u|^3}{3(1-|u|)}
 \leq\frac{\ell^3}{2L^3}.
\]
Also
\[
 0\leq u+\frac\ell L\leq\frac{2\ell}{L^2},\qquad
 \left|\frac{u^2}{2}-\frac{\ell^2}{2L^2}\right|
 \leq\frac{2\ell^2}{L^3}.
\]
Substitution in $r_0+\log r_0-L$ cancels the terms through order
$L^{-2}$ and yields
\[
 |r_0+\log r_0-L|
 \leq\frac{(5/2)\ell^2+(1/2)\ell^3}{L^3}
 \leq\frac{3\ell^3}{L^3}.
\]
The derivative of $x+\log x$ is at least one on $x>0$, so
\begin{equation}\label{eq:W-expansion-error}
 \left|r-\left(L-\ell+\frac\ell L+
             \frac{\ell^2/2-\ell}{L^2}\right)\right|
 \leq\frac{3\ell^3}{L^3}.
\end{equation}

For completeness, the reciprocal term can be controlled without
substituting an unquantified expansion.  Put $v=(L-r)/L$.
By \eqref{eq:growth-r-bracket}, $0\leq v\leq\ell/L\leq1/4$, and
\[
 0\leq r-(L-\ell)=-\log(1-v)
 \leq\frac{4\ell}{3L}.
\]
The exact identity $(1-v)^{-1}=1+v+v^2/(1-v)$ therefore gives
\begin{equation}\label{eq:W-reciprocal-error}
 \left|\frac1r-\left(\frac1L+\frac\ell{L^2}\right)\right|
 \leq\frac{3\ell^2}{L^3}.
\end{equation}
Adding \eqref{eq:W-expansion-error} and
\eqref{eq:W-reciprocal-error} shows that the error in
$r-1+1/r$ is at most $6\ell^3/L^3$; the terms
$-\ell/L^2$ and $+\ell/L^2$ cancel.  Theorem
\ref{thm:log-envelope} proves
\eqref{eq:debruijn-sharp-error}.

Finally $\ell^2\leq L$ for $L\geq16$.  Also $e^L\geq L^5/120$,
so
\[
 \frac{L+4}{e^L}\leq\frac{150}{L^4}
 \leq\frac\ell{L^2}.
\]
This proves \eqref{eq:debruijn-displayed-error}.  The inequalities used
here are elementary monotonicity inequalities on $L\geq16$, with
their initial cases checked by the exponential series.
\end{proof}

\begin{corollary}[Principal logarithmic and root growth]
\label{cor:root-growth}
With explicit convergence rates supplied by
\eqref{eq:debruijn-sharp-error},
\[
 \frac{\log B_n}{n\log n}\longrightarrow1,
 \qquad
 \log\left(\frac{B_n^{1/n}\log n}{n}\right)
 \longrightarrow-1,
 \qquad
 \frac{B_n^{1/n}\log n}{n}\longrightarrow e^{-1}.
\]
In particular, the ordinary power series $\sum B_nx^n$ has zero
radius of convergence.  These logarithmic conclusions alone do not
imply a relative-error asymptotic formula for $B_n$; the next
subsections supply the additional estimate.
\end{corollary}

\begin{proof}
Divide the explicit expression in Theorem \ref{thm:debruijn} by $L$
for the first assertion.  Add $\ell-L$ for the second, and use the
exponential-series error bound for the third.  To see explicitly that
the rates are primitive recursive, one may bound $\log L\leq\sqrt L$
for $L\geq16$ and then choose $L$ larger than a fixed integer power of
the reciprocal requested error.  The threshold $n=2^{2M}$ ensures
$\log n\geq M$ because $\log2\geq1/2$.
For zero radius, the third assertion makes $B_n^{1/n}$ larger than
any prescribed positive rational number beyond an explicitly bounded
index.  Hence, for every nonzero rational $x$, the terms $B_nx^n$
fail the necessary condition of tending to zero.
\end{proof}

\subsection{The Gaussian constant and an effective Stirling estimate}

Two elementary analytic calculations used below are included to make
the normalization and the arithmetic content explicit.  First,
\begin{equation}\label{eq:gaussian-integral}
 \int_{-\infty}^{\infty}e^{-t^2/2}\,dt=\sqrt{2\pi},\qquad
 \int_T^\infty e^{-t^2/2}\,dt\leq\frac{e^{-T^2/2}}T
 \quad(T>0).
\end{equation}
For the tail bound, use $1\leq t/T$ and integrate
$t e^{-t^2/2}$.  For the constant, square the integral truncated to
$[-R,R]$.  The disk of radius $R$ is contained in that square, which
is contained in the disk of radius $\sqrt2R$.  Integration over
annuli, using their explicit areas, gives
\[
 2\pi(1-e^{-R^2/2})
 \leq\left(\int_{-R}^R e^{-t^2/2}\,dt\right)^2
 \leq2\pi.
\]
Both the annular computation and multiplication of the two integrals
are limits of finite Riemann sums with uniform elementary error
bounds.  Taking a predetermined $R$ large enough for the displayed
tail proves \eqref{eq:gaussian-integral}, without a completeness or
choice principle.

\begin{lemma}[Stirling's formula with an explicit remainder]
\label{lem:stirling-effective}
For every integer $m\geq1$ there is a represented number
$\vartheta_m$ with $0\leq\vartheta_m\leq1/(12m)$ such that
\begin{equation}\label{eq:stirling-effective}
 m!=\sqrt{2\pi m}\left(\frac me\right)^m e^{\vartheta_m}.
\end{equation}
\end{lemma}

\begin{proof}
Define
\[
 A_m=\log(m!)-(m+1/2)\log m+m.
\]
For $x=(2m+1)^{-1}$, the geometric series for
$\log((1+x)/(1-x))$ gives
\[
 A_m-A_{m+1}
 =(m+1/2)\log(1+1/m)-1
 =\sum_{j\geq1}\frac{x^{2j}}{2j+1}.
\]
This is positive and at most
\[
 \frac{x^2}{3(1-x^2)}=\frac1{12m(m+1)}.
\]
Consequently the sequence $A_m$ has a computable limit $c$, and
\[
 0\leq A_m-c\leq\frac1{12m}.
\]
This notation for the limit abbreviates the finite estimate
$0\leq A_m-A_M\leq1/(12m)-1/(12M)$ for $M\geq m$.

We identify $c$ by the elementary Wallis calculation.  Put
$I_j=\int_0^{\pi/2}\sin^j t\,dt$.  Integration by parts gives
$jI_j=(j-1)I_{j-2}$, with $I_0=\pi/2$ and $I_1=1$.
Since $0\leq\sin t\leq1$,
\[
 I_{2m+1}\leq I_{2m}\leq I_{2m-1}
 =\frac{2m+1}{2m}I_{2m+1}.
\]
The recurrence also gives the exact identities
\[
 I_{2m}=\frac\pi2\frac{\binom{2m}{m}}{4^m},
 \qquad I_{2m}I_{2m+1}=\frac\pi{2(2m+1)}.
\]
These inequalities, with error at most $1/(2m)$ in the ratio, imply
\[
 \sqrt m\,\frac{\binom{2m}{m}}{4^m}\longrightarrow\frac1{\sqrt\pi}.
\]
On the other hand, direct cancellation gives
\[
 A_{2m}-2A_m
 =\log\left(\sqrt m\,\frac{\binom{2m}{m}}{4^m}\right)
 -\frac12\log2.
\]
The left side tends to $-c$ with error at most $5/(24m)$.
The explicit Wallis inequalities identify the right limit as
$-\tfrac12\log(2\pi)$.  Thus $c=\tfrac12\log(2\pi)$, which proves
\eqref{eq:stirling-effective}.  Every limiting step in this proof has
the displayed reciprocal-integer error bound.
\end{proof}

\subsection{A quantitative saddle-point theorem}

Fourier coefficient extraction from the exponential generating
function gives
\begin{equation}\label{eq:growth-fourier}
 \frac{B_n}{\mathcal A_n}
 =\sqrt{\frac b{2\pi}}\int_{-\pi}^{\pi}e^{\Phi(\theta)}\,d\theta,
 \quad
 \mathcal A_n=\frac{n!\,e^{a-1}}{r^n\sqrt{2\pi b}},
 \quad
 \Phi(\theta)=e^{r e^{i\theta}}-a-in\theta.
\end{equation}
This identity needs no general complex integration theorem.  Expand
the entire function in its explicitly convergent power series on
$|z|=r$ and integrate each finite truncation.  The integral of
$e^{i(j-n)\theta}$ is zero for $j\ne n$ and is $2\pi$ for $j=n$.
The exponential-series tail bounds justify the passage to the
represented integral.  Although the integrand is complex, the
integral is real because the values at $\theta$ and $-\theta$ are
complex conjugates.

\begin{theorem}[A four-term relative-error estimate]
\label{thm:quantitative-saddle}
Suppose $r\geq1$, $T\geq1$, and $T\leq\sqrt b/r$.  Put
\[
 \epsilon=T^3\sqrt{\frac rn}.
\]
Then
\begin{align}\label{eq:saddle-error}
 \left|\frac{B_n}{\mathcal A_n}-1\right|
 \leq{}&e^\epsilon-1+e^{-T^2/2}
       +\frac{24}{T}e^{-T^2/48}
       +3\sqrt b\,e^{-n/(6r^2)}.
\end{align}
\end{theorem}

\begin{proof}
Set $\delta=T/\sqrt b$.  The first two derivatives of the phase at
zero are $\Phi'(0)=0$ and $\Phi''(0)=-b$.  For real $\theta$,
\[
 \left|\frac{d^3}{d\theta^3}e^{r e^{i\theta}}\right|
 \leq e^r(r+3r^2+r^3)\leq5nr^2.
\]
Taylor's formula therefore gives, on $|\theta|\leq\delta$,
\[
 \Phi(\theta)=-\frac{b\theta^2}{2}+E(\theta),
 \qquad |E(\theta)|\leq\frac56nr^2\delta^3\leq\epsilon.
\]
The difference between the normalized integral over this central
interval and the corresponding Gaussian integral is at most
$e^\epsilon-1$.  Completing the Gaussian integral to the whole line
costs at most $e^{-T^2/2}$, by
\eqref{eq:gaussian-integral}.

There are two tail regions.  On $|\theta|\leq1/r$ we have
$|r\sin\theta|\leq1$, $|\sin\theta|\geq|\theta|/2$, and
$e^{r\cos\theta}\geq a/2$.  Since $\cos u\leq1-u^2/3$ for
$|u|\leq1$,
\begin{align}
 \Re\Phi(\theta)
 &=(e^{r\cos\theta}-a)
   +e^{r\cos\theta}(\cos(r\sin\theta)-1)\notag\\
 &\leq-\frac{ar^2\theta^2}{24}
 \leq-\frac{b\theta^2}{48}.
 \label{eq:growth-middle-tail}
\end{align}
Thus the two pieces $\delta\leq|\theta|\leq1/r$ contribute at most
\[
 \frac2{\sqrt{2\pi}}\int_T^\infty e^{-u^2/48}\,du
 \leq\frac{24}{T}e^{-T^2/48}.
\]

For $1/r\leq|\theta|\leq\pi$, use
$1-\cos\theta\geq2\theta^2/\pi^2\geq\theta^2/5$.
Writing $s=r(1-\cos\theta)$, we get $s\geq1/(5r)$, and
\begin{equation}\label{eq:growth-far-tail}
 \Re\Phi(\theta)
 \leq-a(1-e^{-s})
 \leq-\frac a{6r}=-\frac n{6r^2}.
\end{equation}
We used $1-e^{-s}\geq s/(1+s)$, which follows from $e^s\geq1+s$.
The normalized length of the full integration interval is
$\sqrt{2\pi b}<3\sqrt b$.  This proves the final term in
\eqref{eq:saddle-error} and completes the proof.
\end{proof}

\begin{theorem}[A primitive recursive modulus for the saddle formula]
\label{thm:saddle-modulus}
For integers $q\geq1$ and
\begin{equation}\label{eq:saddle-modulus}
 n\geq N(q):=2^{64q^2},
\end{equation}
one has
\begin{equation}\label{eq:saddle-relative-modulus}
 \left|\frac{B_n}{\mathcal A_n}-1\right|\leq\frac5{8q}.
\end{equation}
The threshold is deliberately generous; its role is to give a simple
fully specified arithmetic modulus.
\end{theorem}

\begin{proof}
Use $T=16q$.  At and above the threshold,
\[
 \frac n{\log n}\geq2^{30}q^8,
 \qquad n\geq12(\log n)^3,
 \qquad n\geq24q.
\]
For example, $n/\log n$ is increasing for $n>e$; at $N(q)$ its
value is at least $2^{64q^2-6}/q^2$.  The inequality
$q^{10}\leq2^{10q}$ and $64q^2-36\geq10q$ verify the first
bound.  The second follows similarly from
$q^6\leq2^{6q}$ and $64q^2-18-6q\geq40$.  These are elementary
integer inequalities.

Since $r\leq\log n$, the first bound gives
$\epsilon=T^3\sqrt{r/n}\leq1/(8q)$ and also verifies
$T\leq\sqrt b/r$.  Therefore $e^\epsilon-1\leq1/(4q)$.
The next two terms in \eqref{eq:saddle-error} are each at most
$1/(8q)$: for the third, it suffices to check
$\exp(16q^2/3)\geq12$, already implied by the first three terms of
the exponential series.  Finally $r+1\leq n$, so $\sqrt b\leq n$;
and $n\geq12(\log n)^3$ gives
\[
 3\sqrt b\,e^{-n/(6r^2)}
 \leq3n e^{-2\log n}=\frac3n\leq\frac1{8q}.
\]
Adding the four bounds proves the assertion.
\end{proof}

\begin{corollary}[The Lambert $W$ equivalents on the source pages]
\label{cor:lambert-asymptotic}
As $n$ grows, with primitive recursive convergence moduli,
\begin{align}
 B_n&\sim\frac{n!\,e^{e^{W(n)}-1}}
                  {W(n)^n\sqrt{2\pi n(1+W(n))}},
 \label{eq:precise-saddle-equivalent}\\
 B_n&\sim\frac{n!\,e^{e^{W(n)}-1}}
                  {W(n)^n\sqrt{2\pi W(n)^2e^{W(n)}}},
 \label{eq:odlyzko-equivalent}\\
 B_n&\sim\frac1{\sqrt n}
          \left(\frac n{W(n)}\right)^{n+1/2}
          \exp\left(\frac n{W(n)}-n-1\right).
 \label{eq:wikipedia-equivalent}
\end{align}
For each of the last two expressions, call it $\mathcal C_n$,
the concrete assertion
\begin{equation}\label{eq:lambert-relative-modulus}
 q\geq1,\ n\geq2^{64q^2}\quad\Longrightarrow\quad
 \left|\frac{B_n}{\mathcal C_n}-1\right|\leq\frac1q
\end{equation}
holds.
\end{corollary}

\begin{proof}
The first formula is Theorem \ref{thm:saddle-modulus}.
For the second, its denominator replaces $\sqrt{r+1}$ by
$\sqrt r$, so its ratio with the first approximation is
$\sqrt{(r+1)/r}$.  At the stated threshold,
$L\geq32q^2$, and \eqref{eq:growth-r-bracket} gives
$r\geq L/2\geq16q^2$.  In particular,
\[
 0\leq1-\sqrt{\frac r{r+1}}\leq\frac1{2r}.
\]
For the last formula, substitute Lemma
\ref{lem:stirling-effective} into
\eqref{eq:precise-saddle-equivalent}.  If $R=B_n/\mathcal A_n$,
the exact ratio to its proposed right side is
\[
 R\,e^{\vartheta_n}\sqrt{\frac r{r+1}}.
\]
Here $|R-1|\leq5/(8q)$, $0\leq\vartheta_n\leq1/(12n)$, and
$1-1/(2r)\leq\sqrt{r/(r+1)}\leq1$.  The upper bound follows from
$e^{\vartheta_n}\leq1+1/(6n)$; the lower bound follows by
multiplying the nonnegative lower bounds.  Both differ from one by
less than $1/q$, proving \eqref{eq:lambert-relative-modulus}.
\end{proof}

\subsection{Successive ratios}

The logarithmic convexity established earlier is enough to obtain
successive-ratio growth without differentiating an asymptotic error
term.

\begin{corollary}[A quantitative ratio estimate]\label{cor:bell-ratio}
For every integer $n\geq16$,
\begin{equation}\label{eq:bell-ratio-error}
 \left|\log\frac{B_{n+1}}{B_n}-W(n)\right|
 \leq\frac{4\log n+22}{\sqrt n}.
\end{equation}
Consequently
\[
 \frac{B_{n+1}}{B_n}\sim\frac n{W(n)}\sim\frac n{\log n}.
\]
\end{corollary}

\begin{proof}
Let $F(n)=\log B_n$ and $h=\lfloor\sqrt n\rfloor$; the latter is
the integer square root, a bounded-search primitive recursive
operation.  Logarithmic convexity says that the increments of $F$
are nondecreasing.  Hence
\[
 \frac{F(n)-F(n-h)}h\leq F(n+1)-F(n)
 \leq\frac{F(n+h)-F(n)}h.
\]
Extend the explicitly defined $H(x)=x(W(x)-1+1/W(x))$ to positive
represented arguments.  Direct implicit differentiation of
$W(x)e^{W(x)}=x$ gives
\[
 H'(x)=W(x),\qquad
 0<H''(x)=\frac{W(x)}{x(1+W(x))}\leq\frac1x.
\]
Thus either secant slope of $H$ differs from $W(n)$ by at most
$h/(n-h)\leq2/\sqrt n$.  By
Theorem \ref{thm:log-envelope}, each endpoint error $F(j)-H(j)$,
for $n-h\leq j\leq n+h$, has absolute value at most
$\log n+5$.  Since $h\geq\sqrt n/2$, the error in either secant is
at most $(4\log n+20)/\sqrt n$.  This proves
\eqref{eq:bell-ratio-error}.  Exponentiation and
\eqref{eq:growth-r-bracket} give the two equivalents.
All uses of derivatives can equivalently be read as the displayed
finite increment bounds with integral remainders.
\end{proof}

\subsection{Why the growth proofs remain primitive recursive}

The preceding proof uses a fixed finite collection of analytic
functions on explicit bounded intervals, not a general theory of
real analysis.  Exponentials and trigonometric functions are evaluated
by finite Taylor sums with their elementary tails; logarithms use
the geometric-series integral after rational range reduction;
positive square roots and $W$ use bounded rational approximation.
All argument-size and precision bounds are primitive recursive.

For a particularly transparent integration bound, the phase in
\eqref{eq:growth-fourier} satisfies
$\Re\Phi(\theta)\leq0$ and $|\Phi'(\theta)|\leq2n$.
Therefore the integrand $e^{\Phi(\theta)}$ has Lipschitz constant
$2n$ on $[-\pi,\pi]$.  A uniform Riemann sum with $M$ intervals
has error, after the normalization in
\eqref{eq:growth-fourier}, at most $32n^2/M$.  Thus
$M\geq2^{t+6}n^2$ suffices for error at most $2^{-t}$; approximating
the finitely many endpoints and integrand values consumes a further
explicit portion of the error budget.  Gaussian tails are bounded by
\eqref{eq:gaussian-integral}, and all other tails were bounded in
the proofs above.  Consequently every integral equality or
inequality used here expands into rational finite-sum assertions at
a prescribed precision.

The arbitrary precision parameter is kept free during the finite
proof.  Inductions establish bounds for finite sums, products,
Taylor polynomials, and rational approximation programs; their
invariants are decidable arithmetic inequalities.  After all
represented expressions have been replaced by interval endpoints
and denominators cleared, these are quantifier-free induction
instances in the language with primitive recursive function
symbols.  The moduli in \eqref{eq:saddle-modulus},
\eqref{eq:debruijn-sharp-error}, and
\eqref{eq:bell-ratio-error} prevent any hidden use of an unbounded
search for an asymptotic threshold.  This is the intended iPRA
reading of the human proof; compiling these finite certificates
into the repository's object-language derivation type is a
separate formalization task, specified in the formalization section.

